\subsection{Reading Textual Data}

Reading a text file means asking Python to pull stored bytes from an external file resource, decode those bytes into \texttt{str} objects, and return those strings to the running program. Once the text has been read, it behaves like ordinary Python text. It can be printed, sliced, split, counted, searched, and transformed with normal string methods.

The important boundary is still present: the file stores bytes outside the Python process, while the program receives \texttt{str} objects inside the Python runtime.

\subsubsection{Full-File Loading with \texttt{.read()} and Memory Consumption Boundaries}

The \texttt{.read()} method loads the remaining contents of a text file into one \texttt{str} object. This is convenient for small files because the program receives the whole text at once.

\begin{verbatim}
path = "read_full_demo.txt"

with open(path, "w", encoding="utf-8") as file_obj:
    file_obj.write("Line 1: Nobody expects the Spanish Inquisition!\n")
    file_obj.write("Line 2: The answer is 42.\n")
    file_obj.write("Line 3: D'oh!\n")

with open(path, "r", encoding="utf-8") as file_obj:
    content = file_obj.read()

    print("File object while still open:")
    print(f"file_obj       = {file_obj}")
    print(f"type(file_obj) = {type(file_obj)}")

    print()
    print("Selected names from dir(file_obj):")
    for name in ["read", "readline", "readlines", "close", "seek", "tell"]:
        print(f"{name:12} -> {name in dir(file_obj)}")

print()
print("Text returned by read():")
print(content)

print("Object returned by read():")
print(f"type(content) = {type(content)}")
print(f"len(content)  = {len(content)}")

print()
print("Selected names from dir(content):")
for name in ["splitlines", "upper", "replace", "strip", "count", "find"]:
    print(f"{name:12} -> {name in dir(content)}")

print()
print(f"Number of lines: {len(content.splitlines())}")
print(f"Count of letter e: {content.count('e')}")
\end{verbatim}

The result of \texttt{read()} is one string. The file may contain several lines, but Python still returns one \texttt{str} object containing all of them.

This is simple, but it has a cost. If the file is large, \texttt{read()} attempts to place the entire file content into memory at once. For a small configuration file or a short note, that is fine. For a multi-gigabyte log file, it is a poor access pattern because the program may consume too much memory.

The practical rule is direct: use \texttt{read()} when the file is reasonably small and the program truly needs the complete content at once.

\subsubsection{Line-Based Reading with \texttt{.readline()} and Iteration over File Objects}

The \texttt{.readline()} method reads one line at a time. A line usually includes its final newline character if one exists. When there is nothing left to read, \texttt{readline()} returns the empty string \texttt{""}.

\begin{verbatim}
path = "readline_demo.txt"

with open(path, "w", encoding="utf-8") as file_obj:
    file_obj.write("First: Serenity now!\n")
    file_obj.write("Second: No soup for you!\n")
    file_obj.write("Third: 42 is still suspicious.\n")

with open(path, "r", encoding="utf-8") as file_obj:
    line1 = file_obj.readline()
    line2 = file_obj.readline()
    line3 = file_obj.readline()
    line4 = file_obj.readline()

print("Lines returned by readline():")
print(f"line1 = {line1!r}")
print(f"line2 = {line2!r}")
print(f"line3 = {line3!r}")
print(f"line4 = {line4!r}")

print()
print("Object types:")
print(f"type(line1) = {type(line1)}")
print(f"type(line4) = {type(line4)}")

print()
print("Selected names from dir(line1):")
for name in ["strip", "rstrip", "startswith", "split", "replace"]:
    print(f"{name:12} -> {name in dir(line1)}")

print()
print("Lines without trailing newline characters:")
print(line1.rstrip("\n"))
print(line2.rstrip("\n"))
print(line3.rstrip("\n"))
\end{verbatim}

The representation marker \texttt{!r} is used in the formatted string so the newline character is visible as \texttt{\textbackslash n}. Without that representation form, the newline would simply move the output to the next terminal line.

For many files, direct iteration over the file object is cleaner than repeatedly calling \texttt{readline()}.

\begin{verbatim}
path = "file_iteration_demo.txt"

with open(path, "w", encoding="utf-8") as file_obj:
    file_obj.write("Arthur: It is only a model.\n")
    file_obj.write("Guard: Shh!\n")
    file_obj.write("Narrator: The file object can be iterated line by line.\n")

with open(path, "r", encoding="utf-8") as file_obj:
    print("Iterating over the file object:")

    for line_number, line in enumerate(file_obj, start=1):
        clean_line = line.rstrip("\n")
        print(f"{line_number:02d}: {clean_line}")

        print(f"    type(line) = {type(line)}")
\end{verbatim}

File iteration reads one line at a time instead of materializing the whole file as one string or one list. This makes it the normal pattern for processing large text files sequentially.

The file object acts as an iterable source of strings. Each loop step receives one \texttt{str} object representing one line.

\subsubsection{Batch Line Loading with \texttt{.readlines()} and List Materialization Costs}

The \texttt{.readlines()} method reads all remaining lines and returns them as a \texttt{list} of \texttt{str} objects. This is different from \texttt{read()}, which returns one large string.

\begin{verbatim}
path = "readlines_demo.txt"

with open(path, "w", encoding="utf-8") as file_obj:
    file_obj.write("Line A: To infinity, but not always 42.\n")
    file_obj.write("Line B: D'oh!\n")
    file_obj.write("Line C: Spam and eggs.\n")

with open(path, "r", encoding="utf-8") as file_obj:
    lines = file_obj.readlines()

print("Object returned by readlines():")
print(f"lines       = {lines}")
print(f"type(lines) = {type(lines)}")
print(f"len(lines)  = {len(lines)}")

print()
print("Selected names from dir(lines):")
for name in ["append", "extend", "insert", "pop", "sort", "reverse"]:
    print(f"{name:10} -> {name in dir(lines)}")

print()
print("Inspecting each list element:")
for index, line in enumerate(lines):
    print(f"index {index}: {line!r}")
    print(f"    type(line) = {type(line)}")

print()
print("Clean printed version:")
for line in lines:
    print(line.rstrip("\n"))
\end{verbatim}

The returned object is a list. Each element is a string. Because newline characters are part of the file content, each line string usually still contains its trailing \texttt{\textbackslash n}. Methods such as \texttt{rstrip()} or \texttt{strip()} can remove that when needed.

The cost is materialization. \texttt{readlines()} builds a full list containing all lines. That means the program stores the complete file content and the list structure around it. This is useful when the program genuinely needs random access to the lines, repeated passes, indexing, slicing, or list operations.

For simple one-pass processing, direct file iteration is usually better:

\begin{verbatim}
path = "readlines_demo.txt"

print("One-pass processing without readlines():")

with open(path, "r", encoding="utf-8") as file_obj:
    for line_number, line in enumerate(file_obj, start=1):
        clean_line = line.strip()
        print(f"{line_number:02d} | {clean_line}")
\end{verbatim}

The practical distinction is:

\begin{itemize}
    \item \texttt{read()} returns one full \texttt{str};
    \item \texttt{readline()} returns one \texttt{str} at a time;
    \item file iteration yields one \texttt{str} per loop step;
    \item \texttt{readlines()} returns one \texttt{list} containing many \texttt{str} objects.
\end{itemize}

The access method should match the size and purpose of the file. Small complete documents may be read with \texttt{read()}. Large sequential files should usually be processed by iteration. \texttt{readlines()} should be reserved for cases where a real list of all lines is useful.

